%! Choosing a Sample — Unit 1, Lesson 1.3 — Group Activity
\documentclass[10pt]{article}
\usepackage{saar-article}
\usepackage{saar-boxes}
\newcommand{\TallMath}[1]{$\displaystyle #1\rule[-1.4em]{0pt}{3.2em}$}

\begin{document}

\pageheader{Unit 1, Lesson 1.3}{Group Activity}
% NAMESTRIP: no \namedateperiod here. The student writes their name once, on the
% cover this component is stapled behind. See references/conventions.md.

\vspace{0.15in}

\begin{headlinebox}{frost}
\small
\textbf{Four Ways In --- Cedar Ridge Apartments.} Cedar Ridge has $600$
households in $30$ buildings of $20$. Management wants to know how residents feel
about the parking lot. Buildings $1$--$6$ are studios, buildings $7$--$20$ are
one-bedrooms, and buildings $21$--$30$ are two-bedrooms. Every plan below
surveys $60$ households --- the plans differ only in \emph{who gets chosen}.
\end{headlinebox}

\vspace{0.1in}

\boxguard[30]
\begin{tcolorbox}[colback=white, colframe=black!40, breakable,
    title=\textbf{Tier R --- Remediate},
    fonttitle=\bfseries, arc=2mm, left=3mm, right=3mm, top=2mm, bottom=2mm]
\small
\begin{enumerate}[leftmargin=1.6em, itemsep=6pt, topsep=4pt]

  \item Name the two groups in this study.

        Population: \blank{5.5cm} \hfill Sample size: \blank{1.6cm}

  \item Name the sampling technique in each plan. Choose from \emph{simple
        random}, \emph{stratified}, \emph{systematic}, \emph{cluster}, and
        \emph{convenience}.

        \begin{center}
        \renewcommand{\arraystretch}{1.4}
        \begin{tabularx}{0.96\linewidth}{@{} X p{3cm} @{}}
        \toprule
        \textbf{The plan} & \textbf{Technique} \\
        \midrule
        Number all $600$ households and let a computer draw $60$ numbers.
          & \blank{2.6cm} \\
        Draw $3$ building numbers at random; survey all $20$ households in each.
          & \blank{2.6cm} \\
        Take a random sample from the studios, another from the one-bedrooms, and
        another from the two-bedrooms. & \blank{2.6cm} \\
        Start at a random household near the top of the resident list, then take
        every $10$th one after it. & \blank{2.6cm} \\
        Stand by the mailboxes after work and ask the first $60$ people who walk
        past. & \blank{2.6cm} \\
        \bottomrule
        \end{tabularx}
        \end{center}

  \item For the systematic plan, find the interval $k$.

        \begin{work}
          k &= 600 \div 60 \\
            &= 10
        \end{work}

        The random start lands on household $4$. Write the next three picks.

        \blank{2cm} \qquad \blank{2cm} \qquad \blank{2cm}

  \item For the cluster plan, how many households get surveyed in all?

        \begin{work}
          3 \times 20 &= 60 \text{ households}
        \end{work}

  \item One plan in item 2 did not let chance do the choosing. Which one, and who
        chose the sample instead?
        \writeline
\end{enumerate}
\end{tcolorbox}

\vspace{0.1in}

\boxguard
\begin{tcolorbox}[colback=white, colframe=black!40, breakable,
    title=\textbf{Tier A --- Approaching Proficiency},
    fonttitle=\bfseries, arc=2mm, left=3mm, right=3mm, top=2mm, bottom=2mm]
\small
\begin{enumerate}[leftmargin=1.6em, itemsep=5pt, topsep=3pt]

  \item Management chooses the \textbf{stratified} plan, using unit type as the
        strata. The sample is $60$ of $600$ households, so take $1/10$ of each
        stratum.

        \begin{center}
        \renewcommand{\arraystretch}{1.4}
        \begin{tabularx}{0.9\linewidth}{@{} l c X @{}}
        \toprule
        \textbf{Unit type} & \textbf{Households} & \textbf{Number to survey} \\
        \midrule
        Studio      & 120 & \blank{4cm} \\
        One-bedroom & 280 & \blank{4cm} \\
        Two-bedroom & 200 & \blank{4cm} \\
        \midrule
        \textbf{Total} & \textbf{600} & \blank{4cm} \\
        \bottomrule
        \end{tabularx}
        \end{center}

        Show the work for the \textbf{one-bedroom} row.

        \begin{work}
          \dfrac{60}{600} &= \dfrac{1}{10} \\
          \dfrac{1}{10} \times 280 &= 28 \text{ households}
        \end{work}

  \item Two-bedroom households own more cars than studio households do. Explain
        why that fact makes unit type a good choice of strata here.
        \writelines{2}

  \item In every building, household \#1 is the corner unit --- the only one with
        two assigned parking spaces. The resident list runs building by building,
        $20$ households at a time.

        \textbf{(a)} If management ran the systematic plan with an interval of
        $k = 20$ and a random start of $1$, which households would be picked?
        \writeline

        \textbf{(b)} Why is that a problem for a survey about parking?
        \writeline
\end{enumerate}
\end{tcolorbox}

\vspace{0.1in}

\boxguard
\begin{tcolorbox}[colback=white, colframe=black!40, breakable,
    title=\textbf{Tier E --- Extension},
    fonttitle=\bfseries, arc=2mm, left=3mm, right=3mm, top=2mm, bottom=2mm]
\small
\begin{enumerate}[leftmargin=1.6em, itemsep=5pt, topsep=3pt]

  \item Management keeps a list of \emph{buildings} but has never kept a list of
        \emph{residents}. Of the four probability techniques, which ones can they
        still run, and which become impossible? Explain.
        \writelines{2}

  \item Recall how Cedar Ridge is laid out: buildings $1$--$6$ are studios,
        $7$--$20$ are one-bedrooms, $21$--$30$ are two-bedrooms.

        \textbf{(a)} Suppose the cluster plan happens to draw buildings $2$, $5$,
        and $6$. Whose opinions are in the sample of $60$, and whose are missing
        entirely?
        \writeline

        \textbf{(b)} A cluster sample is supposed to work because each cluster is
        a small mix of the whole population. Explain why Cedar Ridge's buildings
        break that assumption.
        \writelines{2}

  \item All four probability plans survey $60$ households. Say what each one buys
        you and what it costs you.

        \begin{center}
        \renewcommand{\arraystretch}{1.4}
        \begin{tabularx}{0.97\linewidth}{@{} p{2.6cm} X X @{}}
        \toprule
        \textbf{Technique} & \textbf{What it buys} & \textbf{What it costs} \\
        \midrule
        Simple random & \blank{4.6cm} & \blank{4.6cm} \\
        Stratified    & \blank{4.6cm} & \blank{4.6cm} \\
        Systematic    & \blank{4.6cm} & \blank{4.6cm} \\
        Cluster       & \blank{4.6cm} & \blank{4.6cm} \\
        \bottomrule
        \end{tabularx}
        \end{center}
\end{enumerate}
\end{tcolorbox}

\end{document}
